\documentclass[12pt,a4paper]{report}
\usepackage[utf8]{inputenc}
\usepackage{amsmath}
\usepackage{amsfonts}
\usepackage{amssymb}
\usepackage{amsthm}
\usepackage{hyperref}

\usepackage{multicol}
\usepackage{fancyhdr}
\usepackage[inline]{enumitem}
\usepackage{tikz}
\usepackage{tikz-cd}
\usetikzlibrary{calc}
\usetikzlibrary{shapes.geometric}
\usetikzlibrary{positioning}
\usepackage[margin=0.5in]{geometry}
\usepackage{xcolor}

\hypersetup{
    colorlinks=true,
    linkcolor=blue,
    filecolor=magenta,      
    urlcolor=cyan,
    pdftitle={Tensors},
    pdfpagemode=FullScreen,
    }

%\urlstyle{same}

\newcommand{\CLASSNAME}{Math 5050 -- Special Topics: Differential Geometry}
\newcommand{\STUDENTNAME}{Paul Carmody}
\newcommand{\ASSIGNMENT}{Section 13: The Bump Function }
\newcommand{\DUEDATE}{October 1, 2025}
\newcommand{\PROFESSOR}{Professor Berchenko-Kogan}
\newcommand{\SEMESTER}{Fall 2025}
\newcommand{\SCHEDULE}{TBD}
\newcommand{\ROOM}{Remote}

\newcommand{\MMN}{M_{m\times n}}
\newcommand{\FF}{\mathcal{F}}

\pagestyle{fancy}
\fancyhf{}
\chead{ \fancyplain{}{\CLASSNAME} }
%\chead{ \fancyplain{}{\STUDENTNAME} }
\rhead{\thepage}
\input{../MyMacros}

\newcommand{\RED}[1]{\textcolor{red}{#1}}
\newcommand{\BLUE}[1]{\textcolor{blue}{#1}}

\newcommand{\SUPP}{\operatorname{supp}}
\newcommand{\CLOSURE}{\operatorname{closure of}}

\begin{document}

\begin{center}
	\Large{\CLASSNAME -- \SEMESTER} \\
	\large{ w/\PROFESSOR}
\end{center}
\begin{center}
	\STUDENTNAME \\
	\ASSIGNMENT -- \DUEDATE\\
\end{center} 

\textbf{\Large{Notes:}}
\begin{description}
	\item \textbf{General form of bump function}
	
	\begin{align*}
		B(x) &= \frac{h}{1+\BARS{\frac{x-c}{r}}^n}
	\end{align*}Where:\\
	\begin{tabular}{l l}
		$h \in (0,1)$ & height \\
		$c \in [a.b]$ & center of bump in $[a,b]$ \\
		$r$ & widens or spreads the curve out \\
		$n \in \N$ & sharpens the corners
	\end{tabular}
	
	\item \textbf{Manifolds use charts as view ports.}  Always keep in mind that much/most of the work must be done on $\R^n$ through a chart $(U, \phi)$.  That is, translate to $\R^n$ via $\phi$, do the necessary analysis (open, compact, etc.) then translate back through $\phi^{-1}$.
	
	\item \textbf{Exercises}
	
	\textbf{Exercise 13.1 (Bump function supported in an open set).}  Let $q$ be a point and $U$ any neighborhood on $q$ in a manifold. Construct a $C^\infty$ bump function at $q$ supported in $U$.
	
	\RED{\textbf{1st attempt: }Starting with equation 13.2 for $\R^n$ on a open ball
	\begin{align*}
		\sigma(x) &= 1-g\PAREN{\frac{\DBARS{x}r^2-a^2}{b^2-a^2}}
	\end{align*}a bump function on the closed ball $\CONJ{B(0,a)}$ with support on $\CONJ{B(0,b)}$ where $a < b$.  Let 
	\begin{align*}
		b &= \max_{x\in \CONJ{U}} \BARS{q-x} \\
		a &= \max_{x\in U} \BARS{q-x} \\
		\sigma(x) &= 1-g\PAREN{\frac{\DBARS{x-q}r^2-a^2}{b^2-a^2}}
	\end{align*}
	}
	\BLUE{\textbf{What I learned from MathGPT:}
	\begin{enumerate}
		\item \textbf{We must `translate' our exercise into $\R^n$ from $M$.} \\ That is, select a chart $(V, \phi)$ such that $q\in V \subset M$ and $\phi: V \to \phi(V) \in \R^n$ is a diffeomorphism.  Also, $V\subset U$ and $\phi(q) = 0$.\footnote{we have to establish a workspace in $\R^n$ not in $M$.  The diffeomorphism allows us to back and forth between $V$ and $\phi(V)$.}
		\item \textbf{Establish the open set within $V$ using a ball.} \\ $B(0,\epsilon) = B(\phi(q), \epsilon)$ establishes such a ball.
		\item \textbf{Using a standard function for bump construction}.\\  Let 
		\begin{align*}
			g(x) &= \BINDEF{e^{-\frac{1}{x}} & t>0}{0 & t \le 0}
		\end{align*}	Define $\rho(x) = g(\epsilon^2-\DBARS{x})/g(\epsilon^2)$.  This function is smooth equals 1 at the origin and is supported over $B(0,\epsilon)$.
		\item \textbf{Define the Bump Function.}\\
		Let $f: M \to \R$ as
		\begin{align*}
			f(p) &= \BINDEF{\rho(\phi(p)) & p \in V}{0 & p \not \in V}
		\end{align*}This function is smooth, equals 1 at $q$, and its support is in $U$.
	\end{enumerate}	}
	
\end{description}

\HLINE
\begin{enumerate}[label=13.\arabic*]

\item \textbf{Support of a finite sum}

	Prove Lemma 13.5
	
	\textbf{Lemma 13.5} If $\rho_1,\dots,\rho_m$ are real-valued functions on a manifold $M$.  Then,
	\begin{align*}
		\SUPP\PAREN{\sum_{i=1}^m \rho_i} = \bigcup_{i=1}^m \SUPP \rho_i.
	\end{align*}
	
	\BLUE{Let's start with just two, $\rho_1$ and $\rho_2$.  Let $\Phi = \rho_1+\rho_2$.  Let $x \in \SUPP \Phi$ then $x \in \SUPP\{\rho_1 + \rho_2\}$.  Thus $\Phi(x) = \rho_1(x)+\rho_2(x)$.  We have three cases, 
	\begin{description}
		\item Case 1: $x \in \SUPP \rho_1$ which implies $x \in \SUPP \rho_1\cup\SUPP \rho_2$.
		\item Case 2: $x \in \SUPP \rho_2$ which implies $x \in \SUPP \rho_1\cup\SUPP \rho_2$.
		\item Case 3: $x \in \SUPP \rho_1\cap\SUPP \rho_2$ which implies $x \in \SUPP \rho_1\cup\SUPP \rho_2$.
	\end{description}	
	This extends naturally to finite sums of $\rho_i$.
	}
	
\item \textbf{Locally finite family and compact set}

Let $\{A_\alpha\}$ be a locally finite family of subsets of a toplogical space $S$.  Show that every compact set $K$ in $S$ has a neighborhood $W$ that intersects only finitely many of the $A_\alpha$.

\BLUE{Let $K$ be a compact subset contained in $S$.  Therefore, for every open cover of $K$ there exists a finite subcover.  That is, for every $x \in K$ such that $x \in \cup A_\alpha$ there exists a finite number $m_x$ and an index list $I_x$, $|I_x|=m_x$ such that $x \cup_{\alpha\in I} A_{\alpha}$ and a function associated with $x, f_x$, such that $\SUPP f_x \subset \subset \cup_{\alpha\in I} A_{\alpha}$. Therefore, there exists $U_x$ such that $\SUPP f_x \subset U_x$.  Let $W \cup U_x$.  Let 
\begin{align*}
	m_W &= \max_{x\in W} m_x.
\end{align*}This number is finite.
}

\item \textbf{Smooth Urysohn Lemma}

\begin{enumerate}[label=(\alph*)]

	\item Let $A$ and $B$ be two disjoint closed sets in a manifold $M$.  Find a $C^\infty$ function $f$ on $M$ such that $r$ is identically 1 on $A$ and identically 0 on $B$.  (\textit{Hint:} Consider $C^\infty$ partition of unity $\{\rho_{M_A}, \rho_{M_B}\}$ subordinate to the open cover $\{M_A, M_B\}$.  This lemma is needed  in Subsection 29.3.)
	
	\item Let $A$ be a closed subset and $U$ an open subset of a manifold $M$.  Show that there is a $C^\infty$ function $F$ on $M$ such that $f$ is identically 1 on $A$ and $\SUPP f \subset U$.
\end{enumerate}
	
\item \textbf{Support of the pullback of a function}
	
	Let $F: N \to M$ be a $C^\infty$ map of manifolds and $h: M \to \R$ a $C^\infty$ real-valued function.  Prove that $\SUPP F^*h \subset F^{-1}(\SUPP h)$. (\textit{Hint:} First show that $(F^*h)^{-1}(\R^\times)\subset F^{-1}(\SUPP h)$.)
	
	\BLUE{Remember that for all $p\in N$,
	\begin{align*}
		F^*h(p) &= (h\circ F)(p) = h(F(p)) \\
		(F^*h)^{-1} &= (h\circ F)^{-1} = F^{-1}\circ h^{-1}
	\end{align*}Restricting $h^{-1}: \R^\times \to M$ gives us $h^{-1}(\R^\times) = \SUPP h$.  Thus, for any $p \in \R^\times$
	\begin{align*}
		(F^*h)^{-1}(p) &\in F^{-1}(\SUPP h) \\
		(F^*h)^{-1}(\R^\times) &\subset F^{-1}(\SUPP h) 
	\end{align*}Clearly, $\SUPP F^*h = (F^*h)^{-1}(\R^\times)$.
	}
	
\item \textbf{Support of the pullback by projection}
	
	Let $f: M \to \R$ be a $C^\infty$ function on a manifold $M$.  If $N$ is another manifold and $\pi: M \times N \to M$ is the projection onto the first factor, prove that 
	\begin{align*}
		\SUPP(\pi^*f)=(\SUPP f)\times N
	\end{align*}
	
	\BLUE{In this case, for every $p \in M\times N$
		\begin{align*}
			\pi^*f(p) &= (f \circ \pi)(p) = f(\pi(p)) \\
			(f \circ \pi)^{-1} &= \pi^{-1}\circ f^{-1} \\
			\text{since } f^{-1}(\R^\times) &= \SUPP f \\
			\pi^{-1}(\SUPP f) &=(f \circ \pi)^{-1}(\R^\times) \times N\\
			&= \SUPP(f\circ \pi) \times N
		\end{align*}
	}
	
	\item \textbf{Pullback of a partition of unity}
	
	Suppose $\{\rho_\alpha\}$ is a partition of unity on a manifold $M$ subordinate to an open cover $\{U_\alpha\}$ of $M$ and $F :N \to M$ is a $C^\infty$ map. Prove that 
	\begin{enumerate}[label=(\alph*)]
		
		\item the collection of supports $\{\SUPP F^* \rho_\alpha\}$ is locally finite.
		
		\item the collection of functions $\{F^*\rho_\alpha\}$ is a partition of unity on $N$ subordinate to the open cover $\{F^{-1}(U_\alpha)\}$ of $N$
		
	\end{enumerate}
		
		\item \textbf{Closure of a locally finite union}
		
		If $\{A_\alpha\}$ is locally finite collection of subsets ina  topological space, then 
		\begin{align*}
			\CONJ{\bigcup A_\alpha} = \bigcup \CONJ{A_\alpha}.
		\end{align*}where $\CONJ{A}$ denotes the closure of the subset $A$.
		
		\textit{Remark.}  For any colection of subsets $A_\alpha$, one always has 
		\begin{align*}
			\bigcup \CONJ{A_\alpha} \subset \CONJ{\bigcup A_\alpha}.
		\end{align*}However, the reverse inclusion is, in general, not true.  For example, Suppose $A_n$ is the closed interval $[0,1-(1/n)]$ in $\R$.  Then 
		\begin{align*}
			\CONJ{\bigcup_{n=1}^\infty A_n} = \CONJ{[0,1]} = [0,1],
		\end{align*}but
		\begin{align*}
			\bigcup_{n=1}^\infty\CONJ{A_n}
=\bigcup_{n=1}^\infty\SQBRACKET{0,1-\frac{1}{•n}} = [0,1)
		\end{align*}If $\{A_\alpha\}$ is a finite collection, the equality (13.4) is easily shown to be true.


\end{enumerate}

\end{document}
